\section{Exercises}
\label{sec:prop-metatheory-06}

\begin{exercise}
\label{ex:prop-metatheory-01}
Parse
\(\neg P\land Q\to R\lor S\) using the declared precedence and associativity
conventions.  Restore every parenthesis required by the grammar and list all
subformulas.
\end{exercise}

\begin{exercise}
\label{ex:prop-metatheory-02}
Prove by structural induction that, in the syntax tree of every formula, the
number of leaves---occurrences of propositional variables or \(\bottom\)---is
one greater than the number of binary-connective occurrences.
\end{exercise}

\begin{exercise}
\label{ex:prop-metatheory-03}
Construct a countervaluation to
\(P\to Q,Q\entails P\).  Explain why this refutes semantic consequence but
does not by itself inspect any derivation.
\end{exercise}

\begin{exercise}
\label{ex:prop-metatheory-04}
Translate ``The process terminates only if either the input is finite or the
search is bounded, but not both.''  Supply a key and state whether ``or'' is
inclusive before adding the final restriction.
\end{exercise}

\begin{exercise}
\label{ex:prop-metatheory-05}
Find conjunctive and disjunctive normal forms for
\((P\to Q)\land(Q\to R)\).  Verify both by a truth table.
\end{exercise}

\begin{exercise}
\label{ex:prop-metatheory-06}
Give a \(\ND_{\CPL}\)-derivation of
\((P\to R)\to((Q\to R)\to(P\lor Q\to R))\).
Mark every discharged assumption.
\end{exercise}

\begin{exercise}
\label{ex:prop-metatheory-07}
Derive \(\neg(P\land\neg P)\) without using \(\mathsf{RAA}\).  Then
identify the discharging rule used in place of \(\mathsf{RAA}\), and explain
why every rule in the derivation is available in both classical and
intuitionistic natural deduction.
\end{exercise}

\begin{exercise}
\label{ex:prop-metatheory-08}
Give a cut-free \(\LK_{\CPL}\)-derivation of
\(P\to Q, Q\to R\Rightarrow P\to R\).
\end{exercise}

\begin{exercise}
\label{ex:prop-metatheory-09}
Prove that logical equivalence is an equivalence relation and a congruence for
every primitive connective.
\end{exercise}

\begin{exercise}
\label{ex:prop-metatheory-10}
Use the maximal-consistent-set construction to find a countervaluation to a
chosen underivable formula.  Show explicitly how membership in the
set determines the valuation.
\end{exercise}

\begin{exercise}
\label{ex:prop-metatheory-11}
Convert
\[
 (P\lor Q)\land(\neg P\lor R)\land(\neg Q\lor R)\land\neg R
\]
to a clause set and give a resolution refutation.
\end{exercise}

\begin{exercise}
\label{ex:prop-metatheory-12}
Show that deleting a tautological clause preserves the set of satisfying
valuations.  Show by example that deleting a clause that some valuation
satisfies, but that is not tautological, can change the set of valuations
satisfying the clause set.
\end{exercise}

